% TOP_PROBLEM: {"id":30007693,"problem_number":"LOCAL-30007693","title":"Cross-disciplinary discovery","category_id":21,"category":{"id":21,"name":"economics","display_name":"Economics","description":"Open economic theory statements, published research agendas, and proposed scoped specifications; question provenance and historical priority are qualified per record.","slug":"economics","order_index":21,"created_at":"2026-10-08T00:00:00Z"},"status":"research_question","external_url":null,"legacy_ids":[],"published":false,"statement":"Estimate whether AI complements mixed-discipline research teams in producing independently validated discoveries net of coordination costs.","economics":{"original_id":182,"field":"Innovation and AI economics","kind":"empirical","formalization_basis":"adapted_research_question","source_status":"explicit_open","priority_status":"earliest_found","priority_note":"Earliest relevant explicit question or agenda passage located in this audit. No claim of first-ever formulation; earlier versions and later solutions were not exhaustively traced.","earliest_verified_reference":{"authors":["Erik Brynjolfsson","Anton Korinek","Ajay K. Agrawal"],"title":"A Research Agenda for the Economic Implications of Transformative Artificial Intelligence","year":2024,"url":"https://digitaleconomy.stanford.edu/app/uploads/2024/11/ETAI-White-Paper.pdf","doi":null,"locator":"Section 3.2, printed pp. 9–10, local-to-global and cross-disciplinary discovery questions","evidence_summary":"Explicitly asks about cross-disciplinary innovation and AI's role in scientific discovery."},"current_reference":{"authors":["Erik Brynjolfsson","Anton Korinek","Ajay K. Agrawal"],"title":"A Research Agenda for the Economic Implications of Transformative Artificial Intelligence","year":2024,"url":"https://digitaleconomy.stanford.edu/app/uploads/2024/11/ETAI-White-Paper.pdf","doi":null,"locator":"Section 3.2, printed pp. 9–10, local-to-global and cross-disciplinary discovery questions","evidence_summary":"Explicitly asks about cross-disciplinary innovation and AI's role in scientific discovery."},"formalization_note":"The source verifies a broader question or research agenda; the population, estimand, equations and restrictions here are our scoped adaptation. This is not a quoted canonical conjecture, and the audit does not prove contemporary unsolved status for every special case."},"statusQualification":"Published question or research agenda verified at the cited locator. The mathematical specification is adapted for this catalogue and includes additional assumptions; source publication does not establish first-ever priority or certify this exact model remains unsolved. The source verifies a broader question or research agenda; the population, estimand, equations and restrictions here are our scoped adaptation. This is not a quoted canonical conjecture, and the audit does not prove contemporary unsolved status for every special case.","statusReviewedAt":"2026-10-08"}
% FIRST_POSTED: 2026-10-08
% ENTRY_KIND: statement_only
% !TeX program = lualatex
\documentclass[11pt]{article}
\usepackage{fontspec}
\usepackage[a4paper,margin=25mm]{geometry}
\usepackage{amsmath,amssymb,mathtools}
\usepackage{xurl}
\usepackage[hidelinks]{hyperref}
\newcommand{\sref}[2]{\hyperlink{source:#1:#2}{[S#2]}}
\setlength{\emergencystretch}{3em}
\begin{document}
% problemId:30007693
\section*{Cross-disciplinary discovery}\label{30007693}
\subsection{Definitions and mathematical statement}
Take interdisciplinary research projects funded by a defined grant program, with comparable budgets and a three-year evaluation horizon. Let \(A\in\{0,1\}\) indicate access to a fixed AI research system, and \(M\in\{0,1\}\) indicate assignment to a team combining two prespecified disciplines rather than a team within one discipline. For project \(i\), let \(Y_i(a,m)\) be the monetary or preregistered scientific value of independently validated discoveries under assignment \((a,m)\), including zero for projects producing none. Define the cross-disciplinary complementarity
\[
\Delta=E[Y_i(1,1)-Y_i(0,1)-Y_i(1,0)+Y_i(0,0)].
\]
The expectation concerns projects eligible for the program; randomized team composition and AI access would identify this finite-population interaction under consistency and negligible spillovers between assigned teams. Where ideas diffuse, define connected team clusters and randomize at that level instead. Measure disciplinary distance using a fixed pre-intervention taxonomy, and evaluate novelty using blinded reviewers rather than citations generated by the intervention itself. Let \(K(m)\) denote the additional coordination cost of mixed teams. Determine whether the distribution of validated gains and \(\Delta\) justify \(K(1)-K(0)\), and which institutional allocation of expertise supports useful discoveries. Positive interaction in this setting would not establish that AI improves every interdisciplinary project.

\par\textbf{Formulation and scope.} The source verifies a broader question or research agenda; the population, estimand, equations and restrictions here are our scoped adaptation. This is not a quoted canonical conjecture, and the audit does not prove contemporary unsolved status for every special case.

\par\textbf{Open-question provenance.} An explicit question or research agenda was verified at the source locator. This does not by itself certify that every later version remains unresolved \sref{30007693}{1}.

\par\textbf{Historical priority.} Earliest relevant explicit question or agenda passage located in this audit. No claim of first-ever formulation; earlier versions and later solutions were not exhaustively traced.

\subsection{Short English statement}
Estimate whether AI complements mixed-discipline research teams in producing independently validated discoveries net of coordination costs.

\subsection{Sources}
\begin{itemize}\raggedright
\item[\textnormal{[S1]}]\hypertarget{source:30007693:1}{} Erik Brynjolfsson, Anton Korinek, Ajay K. Agrawal. 2024. A Research Agenda for the Economic Implications of Transformative Artificial Intelligence. Section 3.2, printed pp. 9–10, local-to-global and cross-disciplinary discovery questions.
\url{https://digitaleconomy.stanford.edu/app/uploads/2024/11/ETAI-White-Paper.pdf}.
{\small\textit{Earliest verified question/agenda reference; historical priority qualified below. Explicitly asks about cross-disciplinary innovation and AI's role in scientific discovery.}}
\end{itemize}
\end{document}
